\section{Subgroup models and application checks}\label{sec:model}
We verify the arithmetic hypotheses of Theorem~\ref{thm:model-comparison}
for the subgroup restrictions, then evaluate their model correlation.
We first work with a single proper subgroup \(H\le G\), of density
\(\rho\), and write \(f=\1_{C_H}\), \(k_H=[G:H]=1/\rho\).
The construction for the two sides is separate. Put
\[
 L=\log N,\quad y=\exp(L^{2/5}),\quad S=\{p:p\le y\},\quad
 E=\exp(\sqrt L),\quad V_S=\prod_{p\le y}(1-1/p).
\]
Let \(n_S\) denote the \(S\)-supported part of \(n\), and
\(n_{S'}=n/n_S\). Define
\[
 E_+(y)=\prod_{\substack{p\le y\\p\in\mathcal P_H}}(1-1/p),\quad
 E_-(y)=\prod_{\substack{p\le y\\p\nmid L_0,\ p\notin\mathcal P_H}}(1-1/p),
 \quad \alpha_0=\varphi(L_0)/L_0.
\]
Then \(V_S=\alpha_0E_+(y)E_-(y)\) for large \(N\).

\subsection{Fractional poles and centered rough means}
Character orthogonality gives, for \(\Re s>1\),
\begin{equation}\label{eq:euler-branches}
 F_H(s)=\prod_{p\in\mathcal P_H}(1-p^{-s})^{-1}
       =G_H(s)\prod_{\eta\in H^\perp}L(s,\eta)^\rho.
\end{equation}
At a forbidden good prime of coset order \(r>1\), the correction
is \((1-p^{-rs})^{1/r}\); thus \(G_H\) converges absolutely for
\(\Re s>1/2\). The literal principal character carries the bad
Euler factors, proving \eqref{eq:marginal-product}. Every
\(\eta\in H^\perp\) equals one on the support of \(f\); every such
twist therefore has the same pole branch.

Define the rough series and its regular part by
\[
 F_R(s)=F_H(s)\prod_{\substack{p\le y\\p\in\mathcal P_H}}(1-p^{-s})
        =(s-1)^{-\rho}H_R(s).
\]
Here \(H_R(1)=\Gamma(\rho)c_HE_+(y)\). Fixed-order Cauchy estimates
on a disk of radius \(c/\log y\) give
\begin{equation}\label{eq:rough-derivatives}
 |H_R^{(j)}(1)|/j!\ll_j H_R(1)(C\log y)^j.
\end{equation}
For a fixed depth \(K\), define on polynomial scales in \(N\)
\begin{equation}\label{eq:hankel-model}
 M_{R,K}(t)=\sum_{j=0}^K\frac{H_R^{(j)}(1)}{j!}
                 \frac{(\log t)^{\rho-j-1}}{\Gamma(\rho-j)}.
\end{equation}
For \(N^a\le t\le N^C\), with fixed \(0<a<C\), this is
\(c_HE_+(y)(\log t)^{\rho-1}(1+O_K(L^{-3/5}))>0\).
Extend it smoothly and nonnegatively below that range with the same
upper bound. Only the polynomial bulk enters principal means.

\begin{lemma}\label{lem:rough-means}
For fixed \(B,A_0>0\), sufficiently large fixed \(K\) has the
following property. If \(b\) is bounded and periodic of period at
most \(L^B\), and \(W\) is a smooth dyadic weight at a polynomial
scale \(T\), with fixed derivatives bounded by fixed powers of \(L\),
then
\begin{align*}
 &\sum_n W(n)b(n)f(n)\1_{(n,\prod_{p\le y}p)=1}\\
 &\quad-\sum_n W(n)b(n)\1_{(n,\prod_{p\le y}p)=1}
       \left(\sum_{\eta\in H^\perp}\eta(n)\right)
                  \frac{M_{R,K}(n)}{V_S}\ll T L^{-A_0}.
\end{align*}
The same holds with a factor \(n^{iu}\), \(|u|\le L^B\).
\end{lemma}
\begin{proof}
Refine the period to \(Q=\lcm(q,L_0)\) and expand on unit residues
by characters. All primes dividing \(Q\) are at most \(y\).
For \(\psi\bmod Q\), the twisted rough series is
\[
 \prod_{\eta\in H^\perp}L(s,\psi\eta)^\rho\,G_\psi(s)
 \prod_{\substack{p\le y\\p\in\mathcal P_H}}(1-\psi(p)p^{-s}).
\]
If \(\psi\) is inflated from \(H^\perp\), this equals \(F_R\);
otherwise all its character factors are nonprincipal.
Use smooth Mellin inversion at height \(\exp(L^{1/10})\), moving
to \(1-c/\log y\). The classical zero-free region is wider there
at these heights and polylogarithmic conductors. Siegel's theorem
places a possible exceptional real zero outside the strip for large
\(N\), after choosing its fixed exponent depending on \(B\).
The finite Euler product costs \(L^{O(1)}\); the displacement gives
\(\exp(-cL^{3/5})\), dominating polynomial \(L\)-function growth at
that height. Smooth tails are smaller than every fixed logarithmic
power. The threshold can be ineffective.

Hankel integration gives \eqref{eq:hankel-model} for each principal
branch. Its relative Taylor error is \(O_K(L^{-3(K+1)/5})\),
by \eqref{eq:rough-derivatives}.
Choose \(K\) after the derivative and character-count costs.
For the unweighted rough population the series
\(\zeta(s)\prod_{p\le y}(1-p^{-s})\) has principal density \(V_S\),
and only the trivial character contributes. The projection
\(\sum_{\eta\in H^\perp}\eta\) therefore reproduces every principal
branch of the first sum with its exact coefficient. The factor
\(n^{iu}\) changes only the derivative budget and choice of \(K\).
\end{proof}
For \(U\subset H\), set
\begin{equation}\label{eq:corrected-model}
 D_U(n)=k_H\1_{n\bmod L_0\in U}f(n_S)
                       \frac{M_{R,K}(n_{S'})}{V_S},\qquad
 \xi_U(n)=f(n)\1_{n\bmod L_0\in U}-D_U(n).
\end{equation}
Then \(D_U\ge0\), \(D_U\ll L^{-3(1-\rho)/5}\) and
\(|\xi_U|\ll1\). Conditioning on an exact record \(n_S=r\le E\)
leaves residual state \(U/r\), so Lemma~\ref{lem:rough-means}
centers its principal mean. Fixed divisor-growth record coefficients
cost only powers of \(L\), since
\(\sum_{r\ S\text{-supported}}d_j(r)/r\ll L^{O_j(1)}\).
Rankin's bound with exponent \(c/\log y\) gives
\begin{equation}\label{eq:record-tail}
\#\{n\le T:n_S>E\}\ll TL^{O(1)}e^{-cL^{1/10}}
 \quad(N^a\le T\le N^C).
\end{equation}
Indeed the Euler moment is \(L^{O(1)}\), whereas
\(E^{-c/\log y}\) gives the exponential saving. The same argument
permits fixed divisor weights. For \(H=G\), use the exact model
\(D_U=\1_{n\bmod L_0\in U}\) and \(\xi_U=0\).

\subsection{The model correlation}
Use separate constructions for the two subgroups with common \(S\).
\begin{proposition}\label{prop:model-correlation}
For a smooth \(0\le\Phi\le1\) supported away from zero and one,
\begin{align*}
 &\sum_{Ax+By=N}\Phi(Ax/N)D_U(x)D_V(y)\\
 &\quad=c_1c_2\Gamma_N C_{\mathcal E}J_{\mathcal E}(N)
       \frac{N}{AB L^\kappa}
       \left(\int_0^1\Phi(t)\,dt+O(L^{-1/2})\right)+O(N^{1/2}).
\end{align*}
The statement is uniform in target valuations. Endpoint strips of
relative length \(\epsilon\) have upper bound
\(O(\epsilon NJ_{\mathcal E}(N)L^{-\kappa}+N^{1/2})\).
For every interval \(I\subset[0,1]\), the unsmoothed model correlation
has the same leading coefficient multiplied by \(|I|\), with absolute
error \(O(NJ_{\mathcal E}(N)L^{-\kappa-\delta_0})\), where
\(\delta_0=1/(8192\varphi(L_0))\).
\end{proposition}
\begin{proof}
Discard records larger than \(E\) using \eqref{eq:record-tail}.
The rough logarithms on the rest are \(L+O(\sqrt L)\), giving
\[
 \frac{M_{R,K}(n_{S'})}{V_S}
 =\frac{c_H L^{\rho-1}}{\alpha_0E_-(y)}(1+O_K(L^{-1/2})).
\]
Count the remaining small-prime sieve set in the fixed residues.
The fundamental lemma applies at level \(N^{\epsilon_0}\), for
small fixed \(\epsilon_0>0\), with
\(u=\log(N^{\epsilon_0})/\log y\asymp L^{3/5}\).
Its relative error is \(e^{-cu\log u}\), and CRT gives endpoint
error \(O(\nu(d))\) at each squarefree sieving modulus \(d\).
Summing those errors costs \(N^{\epsilon_0}L^{O(1)}\).
This is the standard fixed-dimensional fundamental-lemma application
\cite[Theorem 11.7]{richert}; target-root mergers
decrease its dimension and preserve uniformity.

The positive integer line is \((x,y)=(x_0+Bt,y_0-At)\), with
parameter length \(N/(AB)\). Modulo \(L_0\), this parameter bijects
onto all roots: at each prime at least one coefficient is a unit.
If \(R\) roots have the prescribed states, their raw probability is
\(R/L_0\). Multiplying the normalized marginal projections gives
\[
 \frac{R}{L_0}\frac{L_0c_1}{|H_1|}\frac{L_0c_2}{|H_2|}
       =c_1c_2\Gamma_N.
\]
A prime forbidden on only one side has relative factor one.
A prime forbidden on both sides has two excluded roots, merging
exactly when \(p\mid N\), and relative factor
\[
 \frac{1-(2-\1_{p\mid N})/p}{(1-1/p)^2}.
\]
Their product up to \(y\) is
\(C_{\mathcal E}(y)J_{\mathcal E,y}(N)\). The ordinary tail is
\(O(1/y)\), and the target-prime tail is \(O(L/(y\log y))\).
Partial summation inserts \(\Phi\), costing its total variation.
The same upper sieve on an endpoint strip gives its length times
the sparse density, with a power-smaller endpoint error.
To obtain the uniform interval assertion, put
\(\epsilon=L^{-\delta_0}\) and remove strips of this width near
zero and one. On the rest the rough logarithms are
\(L+O(\sqrt L+\delta_0\log L)\), giving the same relative
\(O(L^{-1/2})\) replacement. The fundamental-lemma count in any
remaining interval has the same CRT endpoint error; partial summation
costs the total variation, bounded independently of the interval.
Restoring the strips costs \(O(\epsilon)\) of the sparse scale.
All other tails are superlogarithmic or power-smaller, proving the claim.
\end{proof}

\subsection{Verification of the comparison hypotheses}
The indicators lie between zero and one, and the model bound after
\eqref{eq:corrected-model} gives \(D_U\le1\) eventually for each
fixed depth. Thus (C1) holds, with a threshold allowed to depend on
that depth. The full-unit-group exact model also lies between zero
and one.

Choose banks of primes \(p\equiv1\pmod{L_0}\), with small banks in
\([\exp(L^{1/8}),\exp(L^{1/4})]\) and big banks in
\([\exp(L^{1/3}),\exp(L^{2/5})]\). Delete primes dividing \(N\).
Every bank prime is a unit at \(A,B\), belongs to both allowed
sets, and preserves the fixed states. Hence
\begin{equation}\label{eq:exact-invariance}
 f_U(pn)=f_U(n),\quad D_U(pn)=D_U(n),\quad \xi_U(pn)=\xi_U(n)
\end{equation}
exactly, including higher powers. The deletion changes reciprocal
mass by at most \(O(L/(P_{\min}\log P_{\min}))\), which is smaller
than every fixed inverse logarithmic power. This verifies (C2).

For (C3), fix an \(S\)-smooth record \(r\le E\). If \(f(r)=0\),
both the indicator and model on that record vanish. Otherwise
\(r\bmod L_0\in H\), and insert the residual state \(U/r\)
into the periodic test of Lemma~\ref{lem:rough-means}. Its smooth
scale is \(Y/r\), so the error is \((Y/r)L^{-D}\), as required.
An additional periodic test of period at most \(L^B\) increases the
period by only a fixed factor. For every required accuracy, the
Hankel remainder works for every sufficiently large fixed depth;
constants and thresholds, rather than earlier graph exponents, may
depend on that depth. This proves (C3) with all surviving character
branches retained.

It remains to check the positive envelope (C4). Write
\(\mu_{\mathcal B}\) for the law \eqref{eq:bank-law}.

\begin{lemma}\label{lem:envelope}
For a nonnegative record function \(F(r,s)\le L^B\), and either
indicator or model on each side, its weighted correlation on the
positive line restricted by \(Ax\in J\subset[0,N]\) is at most
\[
 C\frac{|J|J_{\mathcal E}(N)}{AB L^\kappa}\E_{\mu_{\mathcal B}}F
       +NL^{O(1)}e^{-cL^{1/10}}+O(N^{1/2}).
\]
The constant may depend on the fixed coefficients.
\end{lemma}
\begin{proof}
Discard records \(r,s>E\) using \eqref{eq:record-tail} and bounded
weights. There are at most \(E^2\) retained record pairs. A shared
bank prime makes the stratum empty, because no bank prime divides
\(N\). For coprime records, \(r\mid x,s\mid y\) select one
progression of length \(|J|/(ABrs)\). Exactness excludes two roots
at an unselected bank prime and the remaining variable's root at
a prime dividing \(rs\). Thus the bank probability is exactly
\[
 \frac1{rs}\prod_{\substack{p\in\mathcal B\\p\nmid rs}}(1-2/p)
             \prod_{p\mid rs}(1-1/p)=\mu_{\mathcal B}(r,s).
\]
Forbidden-prime slopes remain units. Apply the upper Selberg sieve
with \(z=N^{1/10}\), using primes up to \(w=z^{1/16}\), which
contains all banks for large \(N\). The squared-weight CRT remainder
is \(O(z^{4+\eta})\), for arbitrarily small fixed \(\eta>0\),
uniformly in the records. After summation it is
\(N^{0.45+o(1)}=O(N^{1/2})\).

For two indicators the nonbank sieve denominator has size at least
a constant times \((\log w)^\kappa/J_{\mathcal E}(N)\).
For model times indicator it has size
\((\log y)^{1-\rho_1}(\log w)^{1-\rho_2}/J_{\mathcal E}(N)\);
the pointwise model bound \(L^{-3(1-\rho_1)/5}\) restores precisely
\(L^{-\kappa}\). For two models the denominator is
\((\log y)^\kappa/J_{\mathcal E}(N)\), and the two model bounds
give the same result. At a prime forbidden on one side only that
side's denominator occurs. Target mergers cost at most the full
factor \(J_{\mathcal E}(N)\). Sum the retained main bounds against
\(F\) to obtain the expectation.
\end{proof}

\begin{proof}[Proof of Theorem~\ref{thm:general}]
The preceding checks verify (C1)--(C4), with
\(\mathcal S_N=NJ_{\mathcal E}(N)/(AB L^\kappa)\).
This scale is between \(NL^{-C}\) and \(NL^C\) for a fixed \(C\),
by \(J_{\mathcal E}\ge1\) and Mertens' upper product bound.
Theorem~\ref{thm:model-comparison} replaces both indicators by their
models with error \(O(\mathcal S_NL^{-\delta_0})\), uniformly in
the interval. Proposition~\ref{prop:model-correlation} evaluates that
model count with exactly the coefficient in the theorem.
If \(\Gamma_N=0\), the residue restrictions give zero identically.
Otherwise its finitely many positive values have a fixed positive
minimum, yielding the relative assertion.
\end{proof}

